\section{Discussion, Recommendations and Limitations}
\label{sec:discussion}

\textbf{A blueprint for road agencies.}
For a new car type the results suggest three steps.
(i) Simulate the car over corridors whose profiles are already surveyed, with randomised vehicles and a sensor model; the ranking of segments will transfer.
(ii) Collect about 5\,km of paired drives over profiled road; one pass per segment sufficed here.
(iii) Train jointly on simulated and paired drives, which beat real-only training on average for every simulator variant and every $k$ we tested; a level offset is slightly more accurate when the simulated level error is common to both road classes, which the paired drives reveal by class and a comparison of real and simulated spectra on the calibration roads anticipates.

\textbf{How to evaluate.}
Because a within-route Spearman correlation is unaffected by calibration, it cannot separate a trained estimator from an untrained statistic; IRI estimators should be judged on held-out routes by absolute error and level per road class.
Spatial dependence matters as well: 5\,km-block intervals of paired differences were \infWidthRatioMin--\infWidthRatioMax{} times as wide as segment-level ones, which would have declared the V0-trained model worse than real training.
On fixed roads, real-trained estimators over-correct for speed (LightGBM's estimate changes by \elRealGbmPooledPct\,\% between 40 and 100\,km/h), presumably because smooth roads are driven fast; simulation-trained LightGBM over-corrects less on urban roads, but this difference is not robust to controlling for within-pass speed variation (Supplement~\ref{S-sec:s_spectra}).

\textbf{Limitations.}
The data cover one car model with one type of telematics unit on four routes in one region, road class and route are confounded, and all intervals describe these routes.
Apart from the comparison of the reference simulator with real training, the variant, calibration and spectral analyses were designed after the test routes had been evaluated; they are exploratory and need confirmation on a new car type.
The simulator is a linear half-car without roll; its vibration differs from the real car's at 2--5\,Hz, above 10\,Hz and in its speed scaling.
We compared three ways of combining simulated and real data with one learner; network fine-tuning or domain adaptation might rank differently.
Transfer to other vehicles was examined only in simulation, where the error on 30 unseen vehicles exceeded that on the training vehicles by \vsThreeRel\,\% with 3 and \vsThirtyRel\,\% with 30 training vehicles (Supplement~\ref{S-sec:s_controlled}).
